\section{Dataset}
We use the CellOT-preprocessed Kang lupuspatients scRNA-seq dataset (\texttt{kang-hvg.h5ad}), containing 28{,}871 cells and 1{,}000 highly-variable genes (HVGs). Each cell is annotated with (i) \texttt{condition} $\in \{\texttt{ctrl},\texttt{stim}\}$, (ii) \texttt{sample\_id} (patient id), and (iii) \texttt{cell\_type} (8 immune cell types). We treat \texttt{ctrl} as the source domain and \texttt{stim} as the target domain (control$\rightarrow$stimulated transport). We report results on the full OOD target test set: stimulated cells from the held-out patient \texttt{sample\_id=101} ($n=1324$).

\section{Metrics}
\paragraph{Accuracy (OOD).}
We report \textbf{cell-type classification accuracy} on the OOD target test set (stimulated cells from patient \texttt{101}).

\paragraph{Distributional alignment (RBF MMD$^2$).}
We compute RBF MMD$^2$ between synthetic transported samples $\tilde{Y}$ and the OOD target test cells $Y_{\text{test}}$, over $\gamma \in \{2,1,0.5,0.1,0.01,0.005\}$, and summarize by the minimum and the mean across $\gamma$ (each side subsampled to at most 2{,}000 points).

\section{Results (tuned setting, seed=0, $n_{\text{ref}}=50$)}
We compare three training regimes for the downstream classifier:
\textbf{Ref-only} (train on $n_{\text{ref}}$ labeled target-reference cells),
\textbf{Synth-only} (train on synthetic target-like labeled samples),
and \textbf{Ref+Synth} (train on the union).
The tuned configuration corresponds to \texttt{configs/cellot\_lupus\_kang\_rectifiedflow\_ref50\_tuned.yaml} and uses more clients, longer Stage1/Stage2 training (30 epochs), full target reference, and more synthesis (5{,}000 samples per client). We also disable the noisy label prior and sample server-side labels uniformly (balanced synthetic training set).

\begin{table}[t]
  \centering
  \begin{tabular}{lcccccc}
    \toprule
    $n_{\text{ref}}$ &
    Acc (Ref-only) &
    Acc (Synth-only) &
    Acc (Ref+Synth) &
    $\Delta$ &
    MMD$_{\min}$ &
    MMD$_{\text{mean}}$ \\
    \midrule
    50 & 0.852 & 0.903 & 0.905 & +0.053 & 0.00126 & 0.00727 \\
    \bottomrule
  \end{tabular}
  \caption{OOD target test performance (stimulated cells from held-out patient \texttt{sample\_id=101}, $n=1324$). $\Delta=\text{Acc(Ref+Synth)}-\text{Acc(Ref-only)}$.}
  \label{tab:kang-ref50-tuned}
\end{table}

With $n_{\text{ref}}=50$, synthetic data improves accuracy from 0.852 (Ref-only) to 0.905 (Ref+Synth).

\section{Structure-preservation rebuttal}
To test whether OT removes donor shift without erasing biological structure, we evaluate two additional
metrics on a label-balanced subset of the held-out stimulated test cells and the synthetic samples.
For each cell type, we sample the same number of target-test and synthetic cells, capped at 150 per
domain and requiring at least 10 cells per label.
We then report:
(i) \textbf{label ASW}, the average silhouette width of the combined target and synthetic samples
with respect to \texttt{cell\_type} labels, and
(ii) \textbf{same-label domain mixing@15}, the fraction of opposite-domain neighbors among the
15 nearest neighbors computed \emph{within each cell type} and averaged across labels.
Higher values indicate better target/synthetic overlap for a fixed biological label.

\paragraph{Auxiliary sanity check.}
The analysis script also reports the mean per-cell-type centroid distance between synthetic samples
and the held-out stimulated target cells in the PCA-50 feature space.
This quantity is used as a secondary check that transport moves each cell type toward the target
distribution, but it is not the primary rebuttal metric.

\begin{figure}[t]
  \centering
  \IfFileExists{plots/kang_structure_preservation.pdf}{
    \includegraphics[width=0.98\linewidth]{plots/kang_structure_preservation.pdf}
  }{
    \IfFileExists{../plots/kang_structure_preservation.pdf}{
      \includegraphics[width=0.98\linewidth]{../plots/kang_structure_preservation.pdf}
    }{}
  }
  \caption{
    \textbf{Kang lupuspatients structure-preservation analysis.}
    Left: per-cell-type centroids in the first two PCA coordinates for raw synthetic samples,
    transported synthetic samples, and the held-out stimulated target test cells.
    Arrows show the raw$\rightarrow$transported$\rightarrow$target movement for each cell type.
    Right: label ASW and same-label domain mixing@15 over seeds.
    Transport should increase target/synthetic mixing while keeping cell-type separation intact.
  }
  \label{fig:kang-structure-preservation}
\end{figure}

\paragraph{Reproducibility.}
\begin{verbatim}
python scripts/evaluate_cell_structure_preservation.py \
  --config configs/cellot_lupus_kang_rectifiedflow_ref50_tuned.yaml \
  --seeds 0,1,2 \
  --output-json plots/kang_structure_preservation.json \
  --plot-output plots/kang_structure_preservation.pdf
\end{verbatim}
